\section{Opposite Categories and the Yoneda Lemma}
\label{section-opposite}

\begin{definition}
\label{definition-opposite}
Given a category $\mathcal{C}$ the {\it opposite category}
$\mathcal{C}^{opp}$ is the category with the same objects
as $\mathcal{C}$ but all morphisms reversed.
\end{definition}

\noindent
In other words
$\Mor_{\mathcal{C}^{opp}}(x, y) = \Mor_\mathcal{C}(y, x)$.
Composition in $\mathcal{C}^{opp}$ is the same as in $\mathcal{C}$
except backwards: if $\phi : y \to z$ and $\psi : x \to y$
are morphisms in $\mathcal{C}^{opp}$, in other words arrows
$z \to y$ and $y \to x$ in $\mathcal{C}$,
then $\phi \circ^{opp} \psi$ is the morphism $x \to z$
of $\mathcal{C}^{opp}$ which corresponds to the composition
$z \to y \to x$ in $\mathcal{C}$.

\begin{definition}
\label{definition-contravariant}
Let $\mathcal{C}$, $\mathcal{S}$ be categories.
A {\it contravariant} functor $F$
from $\mathcal{C}$ to $\mathcal{S}$
is a functor $\mathcal{C}^{opp}\to \mathcal{S}$.
\end{definition}

\noindent
Concretely, a contravariant functor $F$ is given
by a map $F : \Ob(\mathcal{C}) \to
\Ob(\mathcal{S})$ and for every morphism
$\psi : x \to y$ in $\mathcal{C}$ a morphism
$F(\psi) : F(y) \to F(x)$. These should satisfy the property
that, given another morphism
$\phi : y \to z$, we have $F(\phi \circ \psi)
= F(\psi) \circ F(\phi)$ as morphisms $F(z) \to F(x)$.
(Note the reverse of order.)

\begin{definition}
\label{definition-presheaf}
Let $\mathcal{C}$ be a category.
\begin{enumerate}
\item A {\it presheaf of sets on $\mathcal{C}$}
or simply a {\it presheaf} is a contravariant functor
$F$ from $\mathcal{C}$ to $\textit{Sets}$.
\item The category of presheaves is denoted $\textit{PSh}(\mathcal{C})$.
\end{enumerate}
\end{definition}

\noindent
Of course the category of presheaves is a proper class.

\begin{example}
\label{example-hom-functor}
Functor of points.
For any $U\in \Ob(\mathcal{C})$ there is a contravariant
functor
$$
\begin{matrix}
h_U & : & \mathcal{C}
&
\longrightarrow
&
\textit{Sets} \\
& &
X
&
\longmapsto
&
\Mor_\mathcal{C}(X, U)
\end{matrix}
$$
which takes an object $X$ to the set
$\Mor_\mathcal{C}(X, U)$. In other words $h_U$ is a presheaf.
Given a morphism $f : X\to Y$ the corresponding map
$h_U(f) : \Mor_\mathcal{C}(Y, U)\to \Mor_\mathcal{C}(X, U)$
takes $\phi$ to $\phi\circ f$. We will always denote
this presheaf $h_U : \mathcal{C}^{opp} \to \textit{Sets}$.
It is called the {\it representable presheaf} associated to $U$.
If $\mathcal{C}$ is the category of schemes this functor is
sometimes referred to as the
\emph{functor of points} of $U$.
\end{example}

\noindent
Note that given a morphism $\phi : U \to V$ in $\mathcal{C}$ we get a
corresponding natural transformation of functors $h(\phi) : h_U \to h_V$
defined by composing with the morphism $U \to V$. This turns
composition of morphisms in $\mathcal{C}$ into composition of
transformations of functors. In other words we get a functor
$$
h :
\mathcal{C}
\longrightarrow
\textit{PSh}(\mathcal{C})
$$
Note that the target is a ``big'' category, see
Remark \ref{remark-big-categories}. On the other hand,
$h$ is an actual mathematical object (i.e.\ a set), compare Remark
\ref{remark-functor-into-sets}.

\begin{lemma}[Yoneda lemma]
\label{lemma-yoneda}
\begin{reference}
Appeared in some form in \cite{Yoneda-homology}. Used by Grothendieck in a
generalized form in \cite{Gr-II}.
\end{reference}
Let $U, V \in \Ob(\mathcal{C})$.
Given any morphism of functors $s : h_U \to h_V$
there is a unique morphism $\phi : U \to V$
such that $h(\phi) = s$. In other words the
functor $h$ is fully faithful. More generally,
given any contravariant functor $F$ and any object
$U$ of $\mathcal{C}$ we have a natural bijection
$$
\Mor_{\textit{PSh}(\mathcal{C})}(h_U, F) \longrightarrow F(U),
\quad
s \longmapsto s_U(\text{id}_U).
$$
\end{lemma}

\begin{proof}
For the first statement, just take
$\phi = s_U(\text{id}_U) \in \Mor_\mathcal{C}(U, V)$.
For the second statement, given $\xi \in F(U)$ define
$s$ by $s_V : h_U(V) \to F(V)$ by sending the element $f : V \to U$
of $h_U(V) = \Mor_\mathcal{C}(V, U)$ to $F(f)(\xi)$.
\end{proof}

\begin{definition}
\label{definition-representable-functor}
A contravariant functor $F : \mathcal{C}\to \textit{Sets}$ is said
to be {\it representable} if it is isomorphic to the functor of
points $h_U$ for some object $U$ of $\mathcal{C}$.
\end{definition}

\noindent
Let $\mathcal{C}$ be a category and let
$F : \mathcal{C}^{opp} \to \textit{Sets}$ be a representable functor.
Choose an object $U$ of $\mathcal{C}$ and an isomorphism $s : h_U \to F$.
The Yoneda lemma guarantees that the pair $(U, s)$
is unique up to unique isomorphism. The object
$U$ is called an object {\it representing} $F$.
By the Yoneda lemma the transformation $s$ corresponds to a unique
element $\xi \in F(U)$. This element is called the {\it universal object}.
It has the property that for $V \in \Ob(\mathcal{C})$ the map
$$
\Mor_\mathcal{C}(V, U) \longrightarrow F(V),\quad
(f : V \to U) \longmapsto F(f)(\xi)
$$
is a bijection. Thus $\xi$ is universal in the sense that every element
of $F(V)$ is equal to the image of $\xi$ via $F(f)$ for a unique morphism
$f : V \to U$ in $\mathcal{C}$.






